% Файл для \input{...}. Не содержит преамбулы и окружения document.
% Совместим с main.tex из сообщения: использует уже объявленные defbox/examplebox/factbox/algobox/warningbox/remarkbox и пакет asymptote.

\providecommand{\R}{\mathbb{R}}
\providecommand{\Q}{\mathbb{Q}}
\providecommand{\Z}{\mathbb{Z}}
\providecommand{\N}{\mathbb{N}}
\providecommand{\F}{\mathbb{F}}
\providecommand{\Span}{\operatorname{span}}
\providecommand{\Lin}{\operatorname{Lin}}

% Подгрузка служебных определений Asymptote на повторном запуске LaTeX.
% Нужна для корректной вставки подписей из inline-Asymptote в некоторых сборках TeX Live.
\InputIfFileExists{\jobname-1.pre}{}{}
\InputIfFileExists{\jobname-2.pre}{}{}
\InputIfFileExists{\jobname-3.pre}{}{}
\InputIfFileExists{\jobname-4.pre}{}{}
\InputIfFileExists{\jobname-5.pre}{}{}
\InputIfFileExists{\jobname-6.pre}{}{}

\chapter*{Азы высшей и линейной алгебры}

\section*{0. Общий язык: множество и операция}

\begin{defbox}{бинарная операция}
Пусть $X$ --- множество. \emph{Бинарная операция} на $X$ --- это правило, которое каждой паре элементов $a,b\in X$ ставит в соответствие элемент $a*b\in X$.

Очень важно, что результат снова лежит в $X$. Это свойство часто называют \emph{замкнутостью} операции.
\end{defbox}

\begin{examplebox}{операции}
\begin{itemize}
    \item На $\Z$ сложение является операцией: если $a,b\in\Z$, то $a+b\in\Z$.
    \item На $\N=\{1,2,3,\ldots\}$ вычитание не является операцией в этом смысле: $2-5=-3\notin\N$.
    \item На множестве ненулевых рациональных чисел $\Q\setminus\{0\}$ деление является операцией: если $a,b\ne 0$, то $a/b\ne 0$.
\end{itemize}
\end{examplebox}

\begin{remarkbox}{}
Одно и то же множество с разными операциями может иметь совершенно разные свойства. Например, $\Z$ со сложением ведёт себя иначе, чем $\Z$ с умножением.
\end{remarkbox}

\section*{1. Группа}

\begin{defbox}{группа}
\emph{Группа} --- это пара $(G,*)$, где $G$ --- множество, а $*$ --- бинарная операция на $G$, для которой выполнены четыре свойства:
\begin{enumerate}
    \item[0.] \textbf{Замкнутость:} если $a,b\in G$, то $a*b\in G$.
    \item \textbf{Ассоциативность:} $(a*b)*c=a*(b*c)$ для любых $a,b,c\in G$.
    \item \textbf{Нейтральный элемент:} существует элемент $e\in G$ такой, что $e*a=a*e=a$ для любого $a\in G$.
    \item \textbf{Обратный элемент:} для каждого $a\in G$ существует элемент $a^{-1}\in G$ такой, что $a*a^{-1}=a^{-1}*a=e$.
\end{enumerate}
\end{defbox}

\begin{remarkbox}
Главное свойство группы --- обратимость. Можно всегда <<вернуться назад>> после совершенного действия.
\end{remarkbox}

\begin{examplebox}{группы}
\begin{itemize}
    \item $(\Z,+)$ --- группа. Нейтральный элемент равен $0$, обратный к $a$ равен $-a$.
    \item $(\Q\setminus\{0\},\cdot)$ --- группа. Нейтральный элемент равен $1$, обратный к $a$ равен $1/a$.
    \item $(\R,+)$ и $(\mathbb{C},+)$ --- группы.
    \item Множество всех движений плоскости, сохраняющих некоторую фигуру, образует группу относительно композиции движений.
\end{itemize}
\end{examplebox}

\begin{warningbox}{не группы}
\begin{itemize}
    \item $(\N,+)$ не группа: у числа $3$ нет обратного элемента по сложению внутри $\N$.
    \item $(\Z,\cdot)$ не группа: у числа $2$ нет целого обратного по умножению.
    \item $(\Q,\cdot)$ не группа: у $0$ нет обратного по умножению.
\end{itemize}
\end{warningbox}

\subsection*{Группа симметрий квадрата как картинка}

У квадрата есть симметрии: повороты и отражения. Если выполнить одну симметрию, а затем другую, получится снова симметрия квадрата. Поэтому симметрии фигуры естественно образуют группу.

\begin{center}
\begin{asy}
size(8cm);
defaultpen(fontsize(10pt));
pair A=(-1,-1), B=(1,-1), C=(1,1), D=(-1,1), O=(0,0);
draw(A--B--C--D--cycle, linewidth(1.0));
dot(O);
label("$O$", O, SE);
label("$1$", A, SW);
label("$2$", B, SE);
label("$3$", C, NE);
label("$4$", D, NW);
draw(arc(O,0.55,15,100), Arrow(size=6));
label("$r$", (0.35,0.55), NE);
draw((-1.35,0)--(1.35,0), dashed+rgb(0.55,0.55,0.55));
draw((0,-1.35)--(0,1.35), dashed+rgb(0.55,0.55,0.55));
label("$r^4=e$", (0,-1.65));
\end{asy}

\small Поворот $r$ на $90^\circ$ порождает циклическую группу $\{e,r,r^2,r^3\}$, где $r^4=e$.
\end{center}

\begin{remarkbox}
В любой группе нейтральный элемент единственен. Кроме того, у каждого элемента обратный элемент тоже единственен.
\end{remarkbox}

\begin{algobox}{как проверить, что структура является группой}
Чтобы доказать, что $(G,*)$ --- группа, удобно идти по списку:
\begin{enumerate}
    \item Проверь, что операция не выводит из $G$.
    \item Докажи ассоциативность. Иногда она наследуется от известной операции, например от сложения чисел или композиции функций.
    \item Найди нейтральный элемент.
    \item Для произвольного элемента $a\in G$ найди обратный элемент и проверь, что он лежит в $G$.
\end{enumerate}
\end{algobox}

\section*{2. Абелева группа}

\begin{defbox}{абелева группа}
Группа $(G,*)$ называется \emph{абелевой}, или \emph{коммутативной}, если для любых $a,b\in G$ выполнено
\[
    a*b=b*a.
\]
\end{defbox}

\begin{examplebox}{абелевы группы}
\begin{itemize}
    \item $(\Z,+)$, $(\Q,+)$, $(\R,+)$ --- абелевы группы.
    \item $(\Q\setminus\{0\},\cdot)$ и $(\R\setminus\{0\},\cdot)$ --- абелевы группы.
    \item Множество остатков $\Z/n\Z$ со сложением по модулю $n$ --- абелева группа.
\end{itemize}
\end{examplebox}

\begin{warningbox}{не всякая группа абелева}
Композиция преобразований далеко не всегда некоммутативна: порядок действий важен. Между <<выпить, потом закусисть>> и <<закусить, потом выпить>> ощутимая разница!
\end{warningbox}

\section*{3. Кольцо}

\begin{defbox}{кольцо}
\emph{Кольцо} --- это множество $R$ с двумя операциями, обычно называемыми сложением и умножением, для которых выполнены свойства:
\begin{enumerate}
    \item $(R,+)$ --- абелева группа. Нейтральный элемент по сложению обозначается $0$.
    \item Умножение ассоциативно: $(ab)c=a(bc)$.
    \item Выполнены дистрибутивные законы:
    \[
        a(b+c)=ab+ac,\qquad (a+b)c=ac+bc.
    \]
\end{enumerate}
Если в кольце есть элемент $1$ такой, что $1\cdot a=a\cdot 1=a$, то говорят, что это \emph{кольцо с единицей}.
\end{defbox}

\begin{remarkbox}
В разных книгах слово «кольцо» иногда означает немного разные вещи: где-то единица $1$ включается в определение, где-то нет.
\end{remarkbox}

\begin{examplebox}{кольца}
\begin{itemize}
    \item $\Z$, $\Q$, $\R$, $\mathbb{C}$ с обычными сложением и умножением --- кольца с единицей.
    \item $\Z/n\Z$ --- кольцо остатков по модулю $n$.
    \item Многочлены $\R[x]$ образуют кольцо: можно складывать и умножать многочлены.
    \item Квадратные матрицы $n\times n$ над $\R$ образуют кольцо. При $n\ge 2$ умножение матриц некоммутативно.
\end{itemize}
\end{examplebox}

\begin{remarkbox}
    В любом кольце
\[
    0a=a0=0,\qquad (-a)b=-(ab),\qquad (-a)(-b)=ab.
\]
Эти свойства не надо добавлять в определение: они выводятся из аксиом.
\end{remarkbox}

\subsection*{Иерархия колец}

\begin{center}
\begin{asy}
size(9cm);
defaultpen(fontsize(10pt));
draw(circle((0,0),2.4), linewidth(1));
draw(circle((-0.55,0),1.55), linewidth(1));
draw(circle((-0.95,0),0.8), linewidth(1));
label("Кольца", (1.25,1.65));
label("Коммутативные кольца", (-0.25,1.05));
label("Поля", (-1.05,0));
label("$\Z$", (1.25,-0.45));
label("$\Q,\ \R,\ \mathbb{C}$", (-1.0,-0.45));
\end{asy}

\small Поля являются специальными коммутативными кольцами с единицей, где каждый ненулевой элемент обратим.
\end{center}

\section*{4. Поле}

\begin{defbox}{поле}
\emph{Поле} --- это множество $F$ с операциями $+$ и $\cdot$, в котором можно выполнять привычную арифметику: складывать, вычитать, умножать и делить на ненулевой элемент. Формально:
\begin{enumerate}
    \item $(F,+)$ --- абелева группа с нейтральным элементом $0$.
    \item $(F\setminus\{0\},\cdot)$ --- абелева группа с нейтральным элементом $1$.
    \item Выполнена дистрибутивность: $a(b+c)=ab+ac$.
    \item $0\ne 1$.
\end{enumerate}
\end{defbox}

\begin{examplebox}{поля}
\begin{itemize}
    \item $\Q$, $\R$, $\mathbb{C}$ --- поля.
    \item $\Z$ не является полем: у $2$ нет обратного целого числа по умножению.
    \item $\Z_p$ является полем тогда и только тогда, когда $p$ --- простое число.
\end{itemize}
\end{examplebox}

\begin{warningbox}{почему $\Z_6$ не поле}
В $\Z/6\Z$ есть ненулевые элементы $2$ и $3$, но
\[
    2\cdot 3\equiv 0\pmod 6.
\]
Значит, ненулевой элемент $2$ не может иметь обратного: если бы существовал $x$ с $2x\equiv 1\pmod 6$, то после умножения на $3$ получили бы $0\equiv 3\pmod 6$, что невозможно.
\end{warningbox}


\begin{center}
\begin{tabularx}{0.98\textwidth}{>{\bfseries}p{0.22\textwidth}X}
\toprule
Структура & Что в ней можно делать \\
\midrule
Группа & Есть одна операция; можно применять элемент и отменять его обратным. \\
Абелева группа & То же самое, но порядок двух элементов можно менять. \\
Кольцо & Есть сложение, вычитание и умножение; сложение образует абелеву группу. \\
Поле & Есть сложение, вычитание, умножение и деление на любой ненулевой элемент. \\
\bottomrule
\end{tabularx}
\end{center}

\newpage

\section*{5. Линейное пространство}

\begin{defbox}{линейное пространство}
Пусть $F$ --- поле. \emph{Линейное пространство} над $F$ --- это множество $V$, элементы которого называются векторами, если заданы:
\begin{itemize}
    \item сложение векторов: $u+v\in V$;
    \item умножение вектора на скаляр: $\lambda v\in V$, где $\lambda\in F$,
\end{itemize}
и выполнены привычные правила:
\begin{align*}
    u+v&=v+u, & (u+v)+w&=u+(v+w), & v+0&=v, & v+(-v)&=0,\\
    \lambda(u+v)&=\lambda u+\lambda v, & (\lambda+\mu)v&=\lambda v+\mu v, & (\lambda\mu)v&=\lambda(\mu v), & 1v&=v.
\end{align*}
\end{defbox}

\begin{remarkbox}
Векторы --- это не обязательно стрелки на плоскости. Векторами могут быть многочлены, функции, последовательности, матрицы. Важно только, что их можно складывать и умножать на скаляры из поля.
\end{remarkbox}

\begin{examplebox}{линейные пространства}
\begin{itemize}
    \item $\R^2$ и $\R^3$ над полем $\R$.
    \item $F^n$ над полем $F$: наборы $(x_1,\ldots,x_n)$.
    \item Многочлены степени не выше $n$ над $\R$:
    \[
        P_n=\{a_0+a_1x+\cdots+a_nx^n\mid a_i\in\R\}.
    \]
    \item Все функции $f\colon X\to\R$ образуют линейное пространство над $\R$.
\end{itemize}
\end{examplebox}

\begin{warningbox}{скаляры зависят от поля}
Одно и то же множество может быть линейным пространством над разными полями с разной размерностью. Например, $\mathbb{C}$ является одномерным пространством над $\mathbb{C}$ и двумерным пространством над $\R$.
\end{warningbox}

\subsection*{Геометрическая картинка сложения и умножения на число}

\begin{center}
\begin{asy}
size(10cm);
defaultpen(fontsize(10pt));
pair O=(0,0), u=(2.4,0.7), v=(0.9,1.7);
draw((-0.4,0)--(3.9,0), rgb(0.55,0.55,0.55), Arrow(size=5));
draw((0,-0.3)--(0,2.8), rgb(0.55,0.55,0.55), Arrow(size=5));
draw(O--u, blue+linewidth(1.1), Arrow(size=6));
draw(O--v, red+linewidth(1.1), Arrow(size=6));
draw(u--(u+v), dashed+red);
draw(v--(u+v), dashed+blue);
draw(O--(u+v), linewidth(1.2), Arrow(size=6));
label("$u$", u, SE);
label("$v$", v, NW);
label("$u+v$", u+v, NE);
dot(O);
label("$0$", O, SW);
\end{asy}

\small Сложение векторов можно видеть как правило параллелограмма.
\end{center}

\section*{6. Линейная комбинация}

\begin{defbox}{линейная комбинация}
Пусть $v_1,\ldots,v_k\in V$ --- векторы линейного пространства над полем $F$. Вектор
\[
    \lambda_1v_1+\lambda_2v_2+\cdots+\lambda_kv_k,
    \qquad \lambda_i\in F,
\]
называется \emph{линейной комбинацией} векторов $v_1,\ldots,v_k$.
\end{defbox}

\section*{7. Линейная оболочка}

\begin{defbox}{линейная оболочка}
Пусть $S=\{v_1,\ldots,v_k\}\subset V$. \emph{Линейная оболочка} системы $S$ --- это множество всех линейных комбинаций её векторов:
\[
    <v_1, \ldots,v_k>= 
    \Lin(v_1, \ldots,v_k)=
    \Span(v_1,\ldots,v_k)=
    \{\lambda_1v_1+\cdots+\lambda_kv_k\mid \lambda_i\in F\}.
\]

\textit{Примечание: в России принято первое обозначение.}
\end{defbox}

\begin{factbox}{оболочка всегда является подпространством}
Линейная оболочка любой системы векторов является линейным подпространством. Более того, это наименьшее подпространство, содержащее все данные векторы.
\end{factbox}

\noindent\textbf{Почему это правда.} Если взять две линейные комбинации данных векторов и сложить их, снова получится линейная комбинация тех же векторов. Если умножить линейную комбинацию на скаляр, снова получится линейная комбинация.

\begin{center}
\begin{asy}
size(10cm);
defaultpen(fontsize(10pt));
pair O=(0,0), u=(1.4,0.6), v=(-0.4,1.5);
draw((-3.0,-1.285)--(3.0,1.285), rgb(0.55,0.55,0.55)+linewidth(1.0));
draw(O--(2.0*u), blue+linewidth(1.1), Arrow(size=6));
draw(O--(-1.5*u), blue+linewidth(0.9), Arrow(size=5));
label("$\Lin(u)$", (2.5,1.15), NE);
label("$u$", 2.0*u, SE);

pair S=(5.0,0);
draw((S+(-2.2,0))--(S+(2.2,0)), rgb(0.55,0.55,0.55), Arrow(size=5));
draw((S+(0,-1.7))--(S+(0,2.1)), rgb(0.55,0.55,0.55), Arrow(size=5));
draw((S+O)--(S+u), blue+linewidth(1.1), Arrow(size=6));
draw((S+O)--(S+v), red+linewidth(1.1), Arrow(size=6));
draw((S+O)--(S+(1.7*u+0.8*v)), linewidth(1.1), Arrow(size=6));
label("$u$", S+u, SE);
label("$v$", S+v, NW);
label("$\Lin(u,v)=\R^2$", S+(0,-1.95));
\end{asy}

\small Один ненулевой вектор на плоскости порождает прямую; два неколлинеарных вектора порождают всю плоскость.
\end{center}

\begin{examplebox}{оболочки в знакомых пространствах}
\begin{itemize}
    \item В $\R^2$ оболочка одного ненулевого вектора --- прямая через начало координат.
    \item В $\R^3$ оболочка двух неколлинеарных векторов --- плоскость через начало координат.
    \item В пространстве многочленов $P_2$ система $1,x,x^2$ порождает всё пространство $P_2$.
\end{itemize}
\end{examplebox}

\section*{8. Полная система векторов}

\begin{defbox}{полная система}
Система векторов $v_1,\ldots,v_k$ в линейном пространстве $V$ называется \emph{полной}, или \emph{порождающей}, если каждый вектор из $V$ является линейной комбинацией этих векторов. Иначе говоря,
\[
    <v_1,\ldots,v_k>=V.
\]
\end{defbox}

\begin{remarkbox}
Полная система отвечает на вопрос: «Можно ли из этих векторов собрать любой вектор пространства?» Если да, система полна. При этом в ней могут быть лишние векторы!
\end{remarkbox}

\begin{examplebox}{полные и неполные системы}
\begin{itemize}
    \item В $\R^2$ система $(1,0),(0,1)$ полна.
    \item В $\R^2$ система $(1,0),(2,0)$ не полна: её линейная оболочка --- только ось $Ox$.
    \item В $\R^2$ система $(1,0),(0,1),(1,1)$ полна, но содержит лишний вектор, потому что $(1,1)=(1,0)+(0,1)$.
    \item В $P_2$ система $1,x,x^2$ полна, а система $1,x$ не полна: многочлен $x^2$ через неё не выразить.
\end{itemize}
\end{examplebox}

\newpage

\section*{9. Линейная независимость}

\begin{defbox}{линейная независимость}
Система $v_1,\ldots,v_k$ называется \emph{линейно независимой}, если равенство
\[
    \lambda_1v_1+\cdots+\lambda_kv_k=0
\]
возможно только при
\[
    \lambda_1=\lambda_2=\cdots=\lambda_k=0.
\]
Если существует нетривиальная комбинация, равная нулю, то система называется \emph{линейно зависимой}.
\end{defbox}


\begin{factbox}{смысл независимости}
Система линейно зависима тогда и только тогда, когда хотя бы один её вектор выражается через остальные.
\end{factbox}

\begin{examplebox}{зависимость и независимость}
\begin{itemize}
    \item Векторы $(1,0)$ и $(0,1)$ в $\R^2$ независимы.
    \item Векторы $(1,0)$, $(0,1)$ и $(1,1)$ в $\R^2$ зависимы, потому что
    \[
        (1,1)-(1,0)-(0,1)=0.
    \]
    \item Любая система, содержащая нулевой вектор, зависима.
    \item В $\R^2$ любые три вектора зависимы.
\end{itemize}
\end{examplebox}

\section*{10. Базис}

\begin{defbox}{базис}
\emph{Базис} линейного пространства $V$ --- это система векторов, которая одновременно:
\begin{enumerate}
    \item полна, то есть порождает всё пространство $V$;
    \item линейно независима, то есть не содержит лишних векторов.
\end{enumerate}
\end{defbox}

\begin{factbox}{главное свойство базиса}
Система $e_1,\ldots,e_n$ является базисом пространства $V$ тогда и только тогда, когда каждый вектор $v\in V$ единственным образом представляется в виде
\[
    v=x_1e_1+x_2e_2+\cdots+x_ne_n.
\]
Числа $x_1,\ldots,x_n$ называются координатами вектора $v$ в этом базисе.
\end{factbox}

\begin{center}
\begin{asy}
size(10cm);
defaultpen(fontsize(10pt));
pair O=(0,0), e1=(1.2,0), e2=(0,1.2), x=2*e1+1.5*e2;
for(int i=0; i<=3; ++i) {
  draw((1.2*i,0)--(1.2*i,2.4), dashed+rgb(0.75,0.75,0.75));
}
for(int j=0; j<=2; ++j) {
  draw((0,1.2*j)--(3.0,1.2*j), dashed+rgb(0.75,0.75,0.75));
}
draw(O--(3.0*e1), rgb(0.55,0.55,0.55), Arrow(size=5));
draw(O--(2.4*e2), rgb(0.55,0.55,0.55), Arrow(size=5));
draw(O--e1, blue+linewidth(1.2), Arrow(size=6));
draw(O--e2, red+linewidth(1.2), Arrow(size=6));
draw(O--x, linewidth(1.2), Arrow(size=6));
draw(2*e1--x, dashed+red);
draw(1.5*e2--x, dashed+blue);
label("$e_1$", e1, S);
label("$e_2$", e2, W);
label("$v=2e_1+1.5e_2$", x, NE);
label("$0$", O, SW);
\end{asy}

\small Базис задаёт систему координат: вектор описывается коэффициентами при базисных векторах.
\end{center}

\begin{examplebox}{базисы}
\begin{itemize}
    \item В $\R^2$ стандартный базис: $e_1=(1,0)$, $e_2=(0,1)$.
    \item В $\R^3$ стандартный базис: $(1,0,0)$, $(0,1,0)$, $(0,0,1)$.
    \item В $P_2$ базис: $1,x,x^2$.
    \item В $\R^2$ векторы $(1,1)$ и $(1,-1)$ тоже образуют базис: они неколлинеарны.
\end{itemize}
\end{examplebox}

\begin{algobox}{как проверить, что система является базисом в $F^n$}
Если мы знаем, что пространство n-мерно, то достаточно проверить одно из двух:
\begin{itemize}
    \item векторы линейно независимы;
    \item векторы порождают всё пространство.
\end{itemize}
В конечномерном пространстве из $n$ векторов каждое из этих свойств автоматически влечёт другое. На практике приводят матрицу из столбцов к ступенчатому виду и проверяют, что ранг равен $n$.
\end{algobox}

\section*{11. Мини-карта понятий}

\begin{center}
\begin{tabularx}{0.98\textwidth}{>{\bfseries}p{0.23\textwidth}X}
\toprule
Понятие & Короткий смысл \\
\midrule
Группа & Одна операция, есть нейтральный и обратные элементы. \\
Абелева группа & Группа, в которой можно менять порядок элементов: $a*b=b*a$. \\
Кольцо & Две операции: сложение образует абелеву группу, умножение ассоциативно, есть дистрибутивность. \\
Поле & Кольцо, где можно делить на любой ненулевой элемент. \\
Линейное пространство & Векторы можно складывать и умножать на скаляры из поля. \\
Линейная комбинация & Сумма вида $\lambda_1v_1+\cdots+\lambda_kv_k$. \\
Линейная оболочка & Все векторы, которые можно получить линейными комбинациями данной системы. \\
Полная система & Система, оболочка которой равна всему пространству. \\
Базис & Полная и линейно независимая система. \\
\bottomrule
\end{tabularx}
\end{center}